\section{Discussion and limitations}
\subsection{What the results change in evaluation practice}
Four practices follow from the evidence, within the conditions tested. First, an interval that feeds a maintenance rule should be evaluated through the rule: both trigger errors, premature and unrecognised, should be reported beside coverage, with engine-level intervals. Coverage was lost under drift and under bias, yet the lower bound erred on the premature side for one sign of shift and on the late side for the other. Second, the share of out-of-range inputs should not be read as a risk indicator; it missed drift on single-condition data and ordered the perturbation families inconsistently on six-condition data. Third, robustness tests should vary the sign pattern of a shift across channels, because the direction of the trigger error followed it; testing only random-sign offsets would have shown coverage of \NBiasEqT--\NBiasEqL{} on FD002 where a coherent positive shift of equal magnitude gave \JCohPosPICPFDtwo. Fourth, mechanisms should not be selected on a single cost ratio: rankings changed with the ratio in \JRevTotal{} of \JRevGroups{} settings, and under drift they did not change only because no mechanism produced a late error.

\subsection{What cannot be inferred}
The results do not show that real sensor drift resembles a positive ramp on every channel; sensor physics is not modelled. They do not show that deployed systems fail in these ways or give a safety margin: one terminal decision per engine says nothing about first-maintenance time, availability or lifetime cost, and the cost is a binary loss. The output swap does not identify a physical mechanism, nor does it show that the models read a shift as degradation; a monitor separating sensor drift from genuine degradation, which also trends, is a separate problem. The results do not extend to N-CMAPSS, other backbones, other rules or other lead-time families.

\subsection{Limitations}
The data are simulated, and two backbones with fixed hyperparameters were studied. The five models of each backbone use five partitions, so the engine-level intervals are conditional on those models, and initialisation and partition variability are confounded. Perturbation levels are fractions of the training range, not calibrated to any instrument; noise, bias and gain have five realisations each, and the sign patterns of Section~\ref{sec:coherent} are limited to three. Engines with fewer than 30 cycles were excluded. The predicted mean is clipped in some mechanisms and not in others (Section~\ref{sec:framework}); coverage under large mean shifts depends on this, and trigger rates of the heteroscedastic head and CP-norm do not. The waveform controls (FD001 and FD002 only; fixed-end shuffle FD002 only) change end value, energy and order together, so temporal order is not identified as a cause. The engine-level analyses were run once, are exploratory and not pre-registered; rank-reversal counts are descriptive. The literature comparison rests on a bounded search. Reproducibility is layered: released result files regenerate the v1.8 tables; the engine-level analyses rerun inference from released checkpoints and reproduce the released aggregate rates; training and a few intermediate files are not covered.
